% !TeX program = xelatex
% ============================================================================
%  第 6 课　复合作为运算：S₃ 是个群　—— 教师版讲义
%  与 02-课件/第06课-复合作为运算/第06课-复合作为运算.tex 同步
%  记号：e 不动，r 转 120°，r² 转 240°，f 沿过 A 的轴翻，
%        rf 沿过 B 的轴翻，r²f 沿过 C 的轴翻；复合 xy 读作「先做 x，再做 y」
% ============================================================================

\jhlesson{6}{复合作为运算：\texorpdfstring{\(S_3\)}{S3} 是个群}{时长 45 分钟\quad·\quad 教具：正三角形纸片\quad·\quad 本课发单页}

\begin{mubiao}
  \begin{itemize}
    \item 能把「先做一个再做另一个」当成一种运算（复合）。
    \item 能逐条验证这 \(6\) 个动作满足群的四条规则。
    \item 能看懂乘法表，并指出「每一行都是同样 \(6\) 个动作的一次重排」。
  \end{itemize}
\end{mubiao}

\noindent{\small 本课节奏：0—6 回顾｜6—20 把复合当运算｜20—34 验证是群｜34—43 乘法表（发单页，跟做）｜43—45 收尾。全程至少 4 次「用手指回答」。}

\section*{一、回顾（0—6 分钟）}

把上节课那 \(6\) 张两行式再摆一遍（不动、转 \(120^\circ\)、转 \(240^\circ\)、三条轴各翻一次）。

一句话过渡：今天换个角度——把「\textbf{先做这个，再做那个}」本身当成一种运算。

\section*{二、把复合当运算（6—20 分钟）}

先给 \(6\) 个动作起名字，写在黑板左边：

\begin{center}
\begin{tabular}{@{}clc@{}}
  \toprule
  名字 & 动作 & 两行式 \\
  \midrule
  \(e\) & 不动 & \(\begin{pmatrix}A&B&C\\A&B&C\end{pmatrix}\) \\[6pt]
  \(r\) & 顺时针转 \(120^\circ\) & \(\begin{pmatrix}A&B&C\\B&C&A\end{pmatrix}\) \\[6pt]
  \(r^2\) & 顺时针转 \(240^\circ\) & \(\begin{pmatrix}A&B&C\\C&A&B\end{pmatrix}\) \\[6pt]
  \(f\) & 沿过 \(A\) 的轴翻 & \(\begin{pmatrix}A&B&C\\A&C&B\end{pmatrix}\) \\
  \bottomrule
\end{tabular}
\end{center}

\begin{definition}[复合]
  两个动作 \(x\) 和 \(y\)，先做 \(x\)、再做 \(y\)，合起来还是一个动作，记作 \(xy\)。
  这个动作叫做 \(x\) 与 \(y\) 的\textbf{复合}。
\end{definition}

\begin{zhu}
  读法要统一：\(xy\) 读作「先做 \(x\)，再做 \(y\)」。这一步定下来，后面整册书都不再解释。
\end{zhu}

\begin{caozuo}
  拿纸片当场做两遍，两个终点用不同颜色标出来：
  \begin{itemize}
    \item 先 \(r\) 后 \(f\)：算出 \(rf=\begin{pmatrix}A&B&C\\C&B&A\end{pmatrix}\)，是沿过 \(B\) 的轴翻折；
    \item 先 \(f\) 后 \(r\)：算出 \(fr=\begin{pmatrix}A&B&C\\B&A&C\end{pmatrix}\)，是沿过 \(C\) 的轴翻折。
  \end{itemize}
  两个结果并排写在黑板上，让学生自己读出结论。然后板书：
  \[
    \mathbf{rf\neq fr} .
  \]
\end{caozuo}

\begin{zhu}
  这就是第 4 课留的那个例子：\textbf{交换律不是群的必要条件}。
  三角形这一组动作，是它最自然的实例。
\end{zhu}

\begin{caozuo}
  手指回答：\(rf\neq fr\) 说明复合不满足哪一条规则？（交换律）
  「转 \(120^\circ\)」的倒着做是什么？（转 \(240^\circ\)，即 \(r^2\)）
  「沿过 \(A\) 的轴翻折」的倒着做是什么？（还是它自己）
\end{caozuo}

\section*{三、验证是群（20—34 分钟）}

四条规则一条一条对：

\begin{itemize}
  \item \textbf{封闭}：两个动作连着做，结果还是这 \(6\) 个之一。\(rf\) 是沿过 \(B\) 的轴翻，
    \(fr\) 是沿过 \(C\) 的轴翻，都没有跑出去。\(\checkmark\)
  \item \textbf{结合}：\((xy)z\) 和 \(x(yz)\) 都是「先 \(x\)，再 \(y\)，再 \(z\)」，结果一样。\(\checkmark\)
  \item \textbf{单位元}：\(e\)（不动）满足 \(xe=ex=x\)。\(\checkmark\)
  \item \textbf{逆元}：\(r\) 的逆是 \(r^2\)，\(r^2\) 的逆是 \(r\)；三个翻折的逆都是它们自己。\(\checkmark\)
\end{itemize}

四条齐了，所以这 \(6\) 个动作构成一个群，通常记作 \(S_3\)。

\begin{dingli}[三角形动作群]
  正三角形的 \(6\) 个对称动作，连同「先做一个再做另一个」这个运算，
  构成一个群，记作 \(S_3\)。并且它\textbf{不满足交换律}。
\end{dingli}

\section*{四、乘法表（34—43 分钟）}

\begin{caozuo}
  \textbf{先发单页。}在黑板上画一张 \(6\times6\) 的空表，横竖都按
  \(e,r,r^2,f,rf,r^2f\) 排好。告诉全班：格子里的动作是「\textbf{先横排的，再做竖排的}」。
  只填 \(r\) 行和 \(f\) 行，一行一行报答案，你只管往格子里写。
\end{caozuo}

\begin{center}
\begin{tabular}{@{}l|cccccc@{}}
  \toprule
   & \(e\) & \(r\) & \(r^2\) & \(f\) & \(rf\) & \(r^2f\) \\
  \midrule
  先 \(r\) & \(r\) & \(r^2\) & \(e\) & \(rf\) & \(r^2f\) & \(f\) \\
  先 \(f\) & \(f\) & \(r^2f\) & \(rf\) & \(e\) & \(r^2\) & \(r\) \\
  \bottomrule
\end{tabular}
\end{center}

\begin{zhu}
  \textbf{不要填满 \(36\) 格}，会把后面的内容挤掉。乘法表真正值钱的地方是下面这一句观察。
\end{zhu}

把两行念一遍，然后问：这一行里，有没有哪个动作出现两次？有没有哪个动作没出现？

一句总结：\textbf{每一行都是同样 \(6\) 个动作，只是重新排了一遍。} 每一列也一样。

\section*{五、收尾（43—45 分钟）}

从 \(f\) 那一行里挑出第二个格子，读出那个关系：

\[
  fr = r^2 f .
\]

白话是：「先翻后转，等于先\textbf{倒着}转再翻。」因为 \(r^2\) 就是 \(r\) 的倒着做，
这句话也可以写成 \(fr=r^{-1}f\)——下节课要用到。

板书三句话收尾：复合是一种运算；这 \(6\) 个动作构成群 \(S_3\)；\(rf\neq fr\)。

\begin{zhu}
  这是第 4 课那四条规则的\textbf{第一个完整例子}：不是数，是动作，规则照样成立。
  这句话说慢一点，它是整门课的转折点。
\end{zhu}

\begin{xiaojie}
  \begin{itemize}
    \item 复合：\(xy\) 表示「先做 \(x\)，再做 \(y\)」，它本身是一种运算。
    \item \(S_3\)：正三角形的 \(6\) 个对称动作按复合构成一个群，且不交换（\(rf\neq fr\)）。
    \item 乘法表：每一行、每一列都是同样 \(6\) 个动作的一次重排，正是「封闭 \(+\) 逆元」的样子。
    \item 一个关系：\(fr=r^2f\)，也就是「先翻后转 \(=\) 先倒着转再翻」。
  \end{itemize}
\end{xiaojie}

\noindent\textbf{下节课}：在 \(S_3\) 里圈出能自成一体的小圈——阶与子群。